\documentclass[b5paper,11pt]{article}
\usepackage{../notestyle}

% 文档专属：页眉文字
\fancyhead[L]{\fontsize{8pt}{10pt}\selectfont\color{gray!70}408数据结构习题解答}
\fancyfoot[L]{\fontsize{8pt}{10pt}\selectfont\color{gray!70} \thepage}

% 文档信息
\title{408数据结构习题解答}
\author{ClaudeRainer}
% 不显示编写日期
\date{}

\begin{document}
\maketitle

\begin{cquote}
    购置了一套王道论坛组编的《2027年数据结构考研复习指导》因此不需要额外整理PDF材料，自己手动整理所有的复习知识点和不必要的笔记。
    因此对之前编写的笔记进行归档，新一套的文档将会直接以对书本中的题目进行解答并且扩展考题的方式进行整理和学习。
\end{cquote}

\section{1.1.3 数据结构三要素 \space 章节试题精选}
数据结构由逻辑结构、存储结构和基本操作（数据的运算）三要素组成。
逻辑结构是数据元素之间的关系，存储结构是数据元素在计算机中的存储方式，基本操作是对数据元素进行处理的操作。

\hd{01 一个完整的数据结构通常应包含以下哪些要素？ C}
\begin{itemize}
    \item[A] 错误，因为数据结构的三要素是逻辑结构、存储结构和基本操作。
    \item[B] 错误，和A的理由相同，选项中只提出了逻辑结构和物理结构，没有提到基本操作。
    \item[C] 正确，\code{数据的逻辑结构、存储结构以及在其上定义的基本操作。}
    \item[D] 错误，理由同上。
\end{itemize}

\begin{solution}
    P2：数据结构包括三方面，逻辑结构、存储结构和数据的运算。 \\
    逻辑结构是数据元素的逻辑关系，即从逻辑角度对数据的描述方式。
    存储结构是数据结构在计算机中的表示结构，也称为物理结构（顺序存储结构、链式存储结构、索引存储结构，散列存储结构）。
    数据的运算包括数据的定义和实现两个方面。
\end{solution}

\hd{02 下面四种数据结构中，（ ）是非线性数据结构。 A 树}
\begin{itemize}
    \item[A] 正确，树是一个非根结点有且仅有一个父结点，但可以有多个后继子结点的非线性数据结构。
    \item[B] 错误，字符串本质上是线性表的一种特殊形式，属于线性数据结构。
    \item[C] 错误，队列是线性表的一种特殊形式，属于线性数据结构。
    \item[D] 错误，栈是线性表的一种特殊形式，属于线性数据结构。 
\end{itemize}

\hd{03 下列选项中，数据逻辑结构的是（ ）  C 有序表}

选项中的顺序表、哈希表和单链表是数据的存储结构，而有序表是数据的逻辑结构。
有序表表示当前结构中的数据元素按照某种逻辑关系进行排列，可以使用顺序表或者链表实现。

\hd{04 下面关于数据结构的说法中，正确的是（ ） A 数据的逻辑结构独立于其存储结构}

A选项正确，数据的逻辑结构是数据元素之间的关系，而存储结构是数据元素在计算机中的存储方式。逻辑结构和存储结构是相互独立的。
同样可以用顺序表举例，顺序表的逻辑结构是线性表，而存储结构可以是顺序存储结构，也可以是链式存储结构。

B选项“数据的存储结构独立于其逻辑结构”是错误的，数据的存储结构是逻辑结构在计算机中的映射，不能独立存在。

而C选项“数据的逻辑结构唯一决定其存储结构”同样是错误的，理由和选项A相同。

D选项“数据结构仅由其逻辑结构和存储结构决定”是错误的，数据结构还包括数据的运算。

\hd{05 在存储数据时，通常不仅要存储个数据元素的值，还要存储（ ） C 数据元素之间的关系}

不知道数据元素之间的关系，无法确定数据元素之间的逻辑结构，也就无法确定数据结构。

A选项中”数据的操作方法“是数据的运算，和存储数据无关。

B选项”数据元素的类型“对于计算机而言，数据的类型是无意义的，创建一个数据只需要知道需要多大的空间然后进行开辟即可，数据元素的类型和存储数据无关。

D选项”数据的存取方法“这并不是在存储数据过程中实现的，如何存取数据是数据的运算，而不是存储数据。

\hd{对于两种不同的数据结构，逻辑结构或物理结构一定不相同吗？}
不一定，数据结构由三个要素组成：逻辑结构、存储结构和数据的运算。两种不同的数据结构，其逻辑结构与物理结构不一定不相同，
当逻辑结构相同和存储结构相同时，只要运算不同，仍是两种不同的数据结构。

例如顺序表和顺序栈的逻辑结构和存储结构均相同，是线性结构。但顺序表允许在任意位置插入和删除元素，
而顺序栈值允许在一端进行插入和删除操作，两者运算规则不同，因此属于不同的数据结构。

同理，顺序栈，顺序队列，顺序串也都是逻辑结构与存储结构相同，但运算不同而区分的例子。

\hd{举例说明对相同的逻辑结构，同一种运算在不同的存储方式下实现时，其效率不同。}

线性表即可以用顺序表实现，也可以用链表实现。线性表的插入删除、顺序存储需要移动元素O(n)； 
链式存储值需要修改指正O(1)，但需要前提——若给定的是位序，两者都需要定位。
若给定指针，单链表在指针之后插入为O(1)，而顺序表仍然是O(n)。因此在相同的逻辑结构下，存储结构不同，运算效率也不同。

% 分页
\newpage{}

\section{1.2.3 算法和算法评价}

\end{document}